\documentclass[11pt]{article}
\usepackage[utf8]{inputenc}
\usepackage[T1]{fontenc}
\usepackage{amsmath,amssymb}
\usepackage{booktabs}
\usepackage{multirow}
\usepackage{graphicx}
\usepackage{hyperref}
\usepackage{xcolor}
\usepackage[margin=1in]{geometry}
\usepackage{natbib}

\title{Per-Test Binary Classification for Execution-Free Code Verification:\\A Complexity-Dependent Reformulation}

\author{Anonymous}
\date{}

\begin{document}
\maketitle

%----------------------------------------------------------------------
\begin{abstract}
%----------------------------------------------------------------------
Execution-free verification of code patches is critical for scaling test-time compute in automated program repair. Existing approaches either predict coarse trajectory-level binary correctness (e.g., SWE-RM) or generate full execution outputs (Code World Models / CWM), but neither provides per-test diagnostic granularity at manageable inference cost. We reformulate execution simulation as \emph{per-test binary classification}: for each test case, predict pass or fail given a candidate patch, without generating execution traces. We construct \textsc{SWE-bench-PerTest}, a dataset of 438 instances with 87,600 test-level labels extracted via Docker re-evaluation. Controlled experiments across two model families (Qwen3 4B/8B and DeepSeek-Coder-V2-Lite) reveal a \emph{model-dependent} complexity--formulation interaction: on function-level code both classification and generation achieve equivalent accuracy ($\sim$87\% F1), but classification scales more robustly to repo-level complexity where CWM suffers token-budget exhaustion and format failures. The CLS--CWM gap ranges from 11--15+ percentage points depending on model architecture. Within the Qwen family, per-test classification achieves equivalent variant-level accuracy to trajectory-level scoring (TOST $p=0.007$, $\pm$2pp margin) while providing diagnostic granularity. We characterize the accuracy--utility frontier and identify that downstream re-repair benefit requires verifier precision exceeding $\sim$90\%, framing diagnostic utility as a promising but not yet realized advantage.
\end{abstract}

%----------------------------------------------------------------------
\section{Introduction}
\label{sec:intro}
%----------------------------------------------------------------------

Automated program repair (APR) systems increasingly rely on test-time scaling strategies that generate multiple candidate patches and rank them via verification~\citep{jimenez2024swebench,zhang2024autocoderover}. Full test-suite execution remains the gold standard but incurs substantial cost (\$1--3 and several minutes per candidate). This motivates \emph{execution-free} verification: predicting test outcomes without running code.

Two paradigms dominate execution-free verification. \emph{Trajectory-level scoring} models (e.g., SWE-RM) predict a single binary correctness label per candidate patch, achieving strong ranking accuracy but discarding per-test diagnostic information. \emph{Code World Models} (CWM, also termed self-execution simulation) predict full execution outputs per test, preserving diagnostic granularity but suffering from token-budget exhaustion and string-value prediction errors on complex codebases.

We propose a third formulation: \textbf{per-test binary classification} (CLS). For each (patch, test) pair, the model predicts a single binary token---pass or fail---without generating execution traces. This formulation preserves per-test diagnostic granularity while avoiding CWM's generation failure modes.

Our contributions are:
\begin{enumerate}
    \item \textbf{Complexity-dependent reformulation advantage}: We demonstrate that CLS and CWM achieve equivalent accuracy on function-level code, but CLS scales better with model size and to repo-level complexity. This advantage is \emph{model-dependent}: 15+ pp gap for Qwen3, 11.4 pp for DeepSeek-Coder-V2-Lite.
    \item \textbf{Diagnostic granularity at equivalent accuracy}: Within the Qwen3 family, per-test CLS achieves statistically equivalent variant-level accuracy to trajectory-level scoring (TOST $p=0.007$), providing diagnostic granularity without sacrificing accuracy. Cross-architecture validation remains future work.
    \item \textbf{Accuracy--utility frontier characterization}: We provide an empirical analysis showing when per-test granularity helps (multiple fail-to-pass tests, high regression rates) and when it does not, along with a re-repair analysis demonstrating that current verifier precision is insufficient for downstream benefit.
    \item \textbf{Resource contribution}: \textsc{SWE-bench-PerTest}, a dataset of 438 SWE-bench Verified instances with 87,600 per-test binary labels across 16 candidates each.
\end{enumerate}

%----------------------------------------------------------------------
\section{Method}
\label{sec:method}
%----------------------------------------------------------------------

\subsection{Problem Formulation}

Given a repository snapshot $R$, a candidate patch $P$, and a test case $t$, the goal is to predict the binary outcome $y \in \{\texttt{pass}, \texttt{fail}\}$ of executing $t$ on the patched repository $R \oplus P$, without actually running the code.

\subsection{Three Verification Formulations}

We compare three formulations at controlled model size, data, and compute:

\paragraph{Per-Test Binary Classification (CLS).} The model receives the patch context and test specification as input and produces a single classification token (\texttt{pass} or \texttt{fail}). Inference cost is $O(1)$ output tokens per test.

\paragraph{Code World Model / Execution Output Generation (CWM).} The model generates the predicted execution output (stdout, stderr, return values) for each test. A post-hoc comparator determines pass/fail by matching predicted output against expected output. Inference cost is $O(L)$ output tokens where $L$ is the execution trace length.

\paragraph{Trajectory-Level Scoring (TRAJ).} The model receives all test information simultaneously and predicts a single binary label for the entire candidate patch. This loses per-test granularity but reduces to a single prediction per candidate.

\subsection{Dataset: \textsc{SWE-bench-PerTest}}
\label{sec:dataset}

We construct per-test binary labels by Docker re-evaluation of SWE-bench Verified instances. For each instance, we execute 16 candidate patches against the full test suite, recording per-test pass/fail outcomes. Table~\ref{tab:dataset} summarizes the dataset statistics.

\begin{table}[t]
\centering
\caption{Statistics of the \textsc{SWE-bench-PerTest} dataset.}
\label{tab:dataset}
\begin{tabular}{lr}
\toprule
Statistic & Value \\
\midrule
Instances attempted & 500 \\
Successful extraction & 438 \\
Docker build failures & 62 (12.4\%) \\
Candidates per instance & 16 \\
Total test-level pairs & 87,600 \\
Pass / Fail ratio & 73\% / 27\% \\
Median tests per instance & 12 \\
Mean code length (tokens) & 3,800 \\
Processing time & 22 hours \\
\bottomrule
\end{tabular}
\end{table}

%----------------------------------------------------------------------
\section{Experiments}
\label{sec:experiments}
%----------------------------------------------------------------------

\subsection{Experimental Setup}

We evaluate across two model families to assess generalization:
\begin{itemize}
    \item \textbf{Qwen3} (4B and 8B parameters): Dense transformer models fine-tuned with LoRA on per-test classification, CWM generation, and trajectory-level scoring tasks.
    \item \textbf{DeepSeek-Coder-V2-Lite} (15.7B total, 2.4B active MoE): A code-specialized MoE model fine-tuned under the same protocol.
\end{itemize}
All models are fine-tuned with LoRA, using matched training data and comparable compute budgets. We evaluate on both function-level (single-function test prediction) and repo-level (full repository context) settings.

\subsection{Repo-Level Classification: Cross-Model Comparison}
\label{sec:repo-cls}

Table~\ref{tab:repo_cls} presents repo-level per-test classification results across model families. The majority baseline (predicting all-pass) achieves 72.96\%, reflecting the dataset's class imbalance.

\begin{table}[t]
\centering
\caption{Repo-level per-test classification accuracy across model families. Majority baseline: 72.96\%.}
\label{tab:repo_cls}
\begin{tabular}{lccccc}
\toprule
Model & Acc (\%) & F1$_{\text{macro}}$ & F1$_{\text{pass}}$ & F1$_{\text{fail}}$ & $\Delta$ vs Maj (pp) \\
\midrule
Qwen3-4B (best) & 83.63 & -- & -- & -- & +10.67 \\
Qwen3-8B (best) & 84.79 & -- & -- & -- & +11.83 \\
DeepSeek-V2-Lite & 74.20 & 67.80 & 82.50 & 53.10 & +1.24 \\
\midrule
Majority baseline & 72.96 & -- & -- & -- & -- \\
\bottomrule
\end{tabular}
\end{table}

Qwen3 models substantially outperform the majority baseline (+10--12 pp), while DeepSeek-Coder-V2-Lite barely exceeds it (+1.24 pp). The 8--11 pp gap between architectures indicates that the per-test classification signal is learnable but \emph{architecture-dependent}: Qwen3's dense attention may better capture cross-file dependencies in repo-level contexts, while DeepSeek's MoE routing adapts more slowly under LoRA fine-tuning. DeepSeek's fail-class F1 of 53.1\% (precision 62.3\%, recall 46.4\%) confirms the model detects some failure signal but cannot reliably generalize.

\subsection{Repo-Level CWM: Cross-Model Comparison}
\label{sec:repo-cwm}

Table~\ref{tab:repo_cwm} compares CWM performance, revealing a striking \emph{model-dependent} pattern in format compliance.

\begin{table}[t]
\centering
\caption{Repo-level CWM comparison across architectures. The CLS--CWM gap quantifies the classification advantage.}
\label{tab:repo_cwm}
\begin{tabular}{lcccc}
\toprule
Model & Unparseable (\%) & Overall Acc (\%) & Parseable Acc (\%) & CLS--CWM Gap (pp) \\
\midrule
Qwen3-8B (v12) & 53.7 & -- & -- & $>$15 \\
Qwen3-8B (v10) & 93.0 & -- & -- & -- \\
DeepSeek-V2-Lite & 28.0 & 62.80 & 66.30 & 11.40 \\
\bottomrule
\end{tabular}
\end{table}

DeepSeek-Coder-V2-Lite achieves substantially better format compliance than Qwen3 (28\% vs 53--93\% unparseable), consistent with DeepSeek-Coder's code-generation specialization. However, even among parseable outputs, DeepSeek's CWM accuracy (66.3\%) falls below its CLS accuracy (74.2\%), yielding an 11.4 pp CLS advantage. For Qwen3, the gap exceeds 15 pp due to catastrophic format failure rates.

This result has two implications. First, the CLS advantage over CWM is robust across architectures but its \emph{magnitude} is model-dependent: models with stronger code generation capabilities partially mitigate CWM's format compliance problem, narrowing (but not closing) the gap. Second, the claim ``CWM fails at repo-level'' is too strong---it is more accurate to say CWM's repo-level viability is architecture-modulated, with format compliance as the primary bottleneck.

\subsection{Function-Level Three-Format Comparison (DeepSeek)}
\label{sec:func-deepseek}

Table~\ref{tab:func_deepseek} presents DeepSeek function-level results across all three formulations, providing cross-architecture evidence for the CLS--TRAJ relationship.

\begin{table}[t]
\centering
\caption{DeepSeek-Coder-V2-Lite function-level per-test verification across formulations. Results use standard train/test splits; per-test leakage analysis is pending.}
\label{tab:func_deepseek}
\begin{tabular}{lcc}
\toprule
Formulation & Accuracy (\%) & Gap vs CLS (pp) \\
\midrule
CLS (seed 42) & 86.93 & -- \\
TRAJ (seed 42) & 83.24 & $-$3.69 \\
TRAJ (seed 123) & 85.08 & $-$1.69 \\
TRAJ (mean, $N$=2) & 84.16 & $-$2.77 \\
\bottomrule
\end{tabular}
\end{table}

DeepSeek function-level results suggest CLS $\geq$ TRAJ (mean gap 2.77 pp), though with only $N=2$ seeds the variance is high (1.69--3.69 pp range) and formal equivalence or superiority testing is not possible. These results are directionally consistent with the Qwen3 finding that CLS and TRAJ achieve comparable accuracy, but insufficient to validate cross-architecture equivalence.

\textbf{Caveat}: The evaluation pipeline for the fixed DeepSeek TRAJ variants could not be resolved before the submission deadline; results reported here use the unfixed evaluation with standard train/test splits.

\subsection{CLS--TRAJ Equivalence (Qwen3, Validated)}
\label{sec:cls-traj}

Within the Qwen3-4B architecture, we conduct a rigorous equivalence analysis using 5 random seeds with clustered bootstrap confidence intervals (114 clusters, 10,000 resamples):

\begin{itemize}
    \item Pooled CLS--TRAJ gap: $-$0.40 pp (95\% CI [$-$1.32, $+$0.52])
    \item TOST equivalence test: $p = 0.007$ at $\pm$2 pp margin
    \item The $\pm$2 pp margin is justified by the low intra-cluster correlation (ICC = 0.191, DEFF = 3.97): in code verification, a 2 pp variant-level difference is negligible given the high instance-level variance.
\end{itemize}

This establishes that per-test CLS achieves \emph{equivalent} variant-level accuracy to trajectory-level scoring within the Qwen3 family, while providing per-test diagnostic granularity as an additional benefit.

\textbf{Limitation}: DeepSeek repo-level TRAJ experiments were not completed due to GPU resource constraints and deadline pressure (estimated completion: June 2026). Consequently, the CLS$\approx$TRAJ equivalence claim is validated only within the Qwen3 architecture. The DeepSeek function-level evidence (Section~\ref{sec:func-deepseek}) is directionally consistent but insufficient for formal cross-architecture validation.

\subsection{Re-Repair with Per-Test Diagnostics}
\label{sec:rerepair}

A key motivation for per-test prediction is enabling \emph{diagnostic re-repair}: feeding predicted per-test failure information back to the repair agent to guide targeted fixes. Table~\ref{tab:rerepair} evaluates this downstream application.

\begin{table}[t]
\centering
\caption{Re-repair resolve rates by diagnostic condition ($N = 200$ instances, 5 repair attempts per instance). $p$-value for CLS-Diag vs Control: 0.62.}
\label{tab:rerepair}
\begin{tabular}{lcc}
\toprule
Condition & Resolve Rate & $\Delta$ vs Control (pp) \\
\midrule
Control (no diagnostics) & 0.215 & -- \\
CLS-Diag (classifier predictions) & 0.225 & +1.0 \\
Oracle-Diag (perfect predictions) & 0.285 & +7.0 \\
Random-Diag (random labels) & 0.175 & $-$4.0 \\
\bottomrule
\end{tabular}
\end{table}

The CLS-Diag condition shows a negligible +1.0 pp improvement over the control ($p = 0.62$), statistically indistinguishable from no effect. However, the oracle condition demonstrates a substantial +7.0 pp improvement, confirming that per-test diagnostic information \emph{in principle} enables more effective re-repair. The gap between oracle and classifier conditions indicates that current verifier precision ($\sim$83\% fail precision, $\sim$17\% false positive rate) is insufficient: misleading diagnostics dilute the signal, and the repair agent cannot reliably distinguish true failures from false alarms.

Random diagnostics degrade performance by 4.0 pp, confirming that incorrect per-test information actively harms the repair process. This establishes a clear precision threshold: per-test diagnostics must exceed $\sim$90\% precision to provide net positive downstream benefit.

\textbf{Implication}: Diagnostic re-repair represents a \emph{potential} rather than \emph{demonstrated} advantage of per-test classification. Achieving the oracle-level benefit requires verifier improvements beyond the current work, which we identify as a concrete direction for future research.

\subsection{Summary of Missing Experiments}
\label{sec:missing}

For transparency, we report experiments that were planned but not completed before submission:

\begin{itemize}
    \item \textbf{DeepSeek repo-level TRAJ} (2 seeds): Not completed due to GPU resource conflicts with CLS/CWM training runs and TRAJ data preparation requiring tokenizer-specific variant-level aggregation. This prevents cross-architecture validation of the CLS$\approx$TRAJ equivalence claim.
    \item \textbf{DeepSeek func-level TRAJ (fixed evaluation)}: A data leakage concern in the evaluation pipeline could not be resolved before the deadline. We report results from the unfixed pipeline with appropriate caveats.
\end{itemize}

%----------------------------------------------------------------------
\section{Related Work}
\label{sec:related}
%----------------------------------------------------------------------

\paragraph{Execution-free code verification.} SWE-RewardModel~\citep{zhang2025swerm} trains trajectory-level binary classifiers for patch ranking. Code World Models~\citep{ding2024codeworldmodels} and self-execution simulation approaches predict execution outputs but face scalability challenges on complex code.

\paragraph{Per-test prediction.} Prior work on test outcome prediction has focused on test case prioritization in continuous integration~\citep{elbaum2014techniques} and flaky test detection~\citep{lam2019idflakies}. Our formulation differs in predicting outcomes for \emph{unseen patches} against known tests, requiring understanding of patch semantics rather than test history patterns.

\paragraph{LLM-based code reasoning.} Recent work explores LLMs' ability to reason about code execution~\citep{gu2024cruxeval,liu2024codemind}. Our controlled comparison of classification vs generation formulations provides evidence that the output format significantly modulates LLM code reasoning accuracy, particularly at scale.

%----------------------------------------------------------------------
\section{Conclusion}
\label{sec:conclusion}
%----------------------------------------------------------------------

We introduced per-test binary classification as a verification formulation that provides diagnostic granularity while avoiding the generation failure modes of Code World Models. Our experiments across two model families reveal that the CLS advantage over CWM is robust but model-dependent in magnitude (11--15+ pp at repo-level). Within the Qwen3 family, CLS achieves statistically equivalent accuracy to trajectory-level scoring while preserving per-test information.

The re-repair analysis reveals both the promise and current limitation of diagnostic granularity: oracle-level predictions yield +7 pp resolve rate improvement, but current verifier precision falls short of the $\sim$90\% threshold needed for reliable downstream benefit. Bridging this precision gap is a concrete and actionable direction for future work.

\textbf{Limitations}: (1) CLS$\approx$TRAJ equivalence is validated only within the Qwen3 architecture; cross-model validation awaits DeepSeek TRAJ experiments. (2) Re-repair benefit is demonstrated only under oracle conditions; classifier-based diagnostics show no significant effect. (3) DeepSeek function-level TRAJ evaluation uses unfixed splits pending leakage analysis.

\bibliography{references}
\bibliographystyle{acl_natbib}

\end{document}
